\clearpage
\section{中斷：從 RTL 的一根線到 driver 的 ISR}
\label{sec:irq}

\begin{oneline}
v1.0 的 bitstream 沒有中斷線，driver 只能每 13--50 µs 讀一次 done 位元。v1.1 在控制暫存器模組 \code{twn_csr} 加了一個 level 中斷輸出，一路接到 PS 的 \code{pl_ps_irq0}（GIC 的 SPI 89）；driver 改成睡在一個 completion 上，由中斷處理函式（ISR）叫醒。
\end{oneline}

\begin{before}
\textbf{硬體}：v1.0.0 的 block design 把 PS 的 PL 中斷輸入接成常數 0，只為了讓 PS 的設定不變：
\snip[tcl]{v10_irq_zero}
\textbf{driver v1.0}：等 done 只能輪詢（第一冊第 6 節）。\code{wait_bit()} 用 \code{readl_poll_timeout()}，每次讀完睡 13 到 50 µs：
\snip{v10_wait_bit}
\snip{v10_wait_done}
問題是一幀 4.1 ms 裡，CPU 要醒來讀暫存器上百次；而且 done 之後最多要再等一個睡眠週期才發現。
\end{before}

\begin{tech}
\begin{itemize}
\item \textbf{RTL}：新暫存器 IRQ（offset \code{0x0064}）放兩個致能位元；中斷輸出是暫存器輸出（registered）、level high。
\item \textbf{IP 封裝}：\code{ipx::infer_bus_interface} 把 \code{interrupt} 這個 port 宣告成 Vivado 的 interrupt 介面，參數 \code{SENSITIVITY = LEVEL_HIGH}。
\item \textbf{block design}：\code{connect_bd_net} 接到 \code{ps/pl_ps_irq0}。PetaLinux 讀 XSA 後自動在 device tree 節點寫 \code{interrupts = <0 89 4>}。
\item \textbf{kernel}：\code{platform_get_irq_optional()}、\code{devm_request_irq()}、ISR 回傳 \code{irqreturn_t}；\code{struct completion} 的 \code{init_completion()}、\code{reinit_completion()}、\code{complete()}、\code{wait_for_completion_timeout()}；module parameter \code{use_irq}。
\end{itemize}
\end{tech}

\begin{gain}
\begin{itemize}
\item \textbf{CPU 時間}：每幀從 \pollCpuUs{} µs 降到 \irqCpuUs{} µs（\code{twn_bench}，同一次開機、同一個程式，只差 \code{use_irq}）。
\item \textbf{等待延遲}：driver 量的時間減掉 PL 自己數的 cycle，從 \pollWaitUs{} µs 變成 \irqWaitUs{} µs。
\item \textbf{吞吐量不變差}：pipeline \fpsIrq{} fps（第 10 節的 board check）。
\item \textbf{相容}：device tree 沒有 \code{interrupts} 的舊 bitstream，driver 自動退回輪詢（開機組 D 的 dmesg 寫 \code{polling}）。
\end{itemize}
\end{gain}

\subsection{硬體：RTL 加一根線}

\code{twn_csr.v} 開頭的說明寫了整個設計：

\snip[Verilog]{csr_irq_doc}

\textbf{中斷怎麼產生。}每個 clock 都把「致能 AND done」重新算一次，存進暫存器 \code{irq}：

\snip[Verilog]{csr_irq_reg}

\code{ap_done} 是 backbone 做完、\code{hd_done_l} 是 head 做完的黏著位元（設起來後一直保持，直到被清掉）。所以 \code{irq} 在 done 設起來的下一個 cycle 拉高，在 done 被清掉的下一個 cycle 放下。中間的 \code{ifdef CSR_IRQ_BREAK_MASK} 是故意做壞的版本，只給模擬用，見下面的 gate。

\textbf{致能放在自己的暫存器。}寫 IRQ 暫存器設定兩個致能，讀它可以同時看到致能、兩個 done 和 \code{irq} 本身：

\snip[Verilog]{csr_irq_wr}
\snip[Verilog]{csr_irq_rd}

為什麼不沿用 CTRL 的 bit 2（\code{irq_en}）？因為 CTRL 的每一次寫入都會把 bit 2 一起寫掉：

\snip[Verilog]{csr_ctrl_wr}

driver 啟動一幀寫的是 \code{0x1}（START），清 done 寫的是 \code{0x2}；這兩個值的 bit 2 都是 0。若致能放在 CTRL[2]，每次啟動都會把中斷關掉。\code{CSR_IRQ_BREAK_CTRL} 就是這個錯誤的版本。

\textbf{單元 gate。}\code{rtl/fuse/sim/run_csr_irq.ps1} 跑三次：原設計必須 PASS，兩個故意做壞的版本必須 FAIL。後兩者證明 testbench 真的能抓到這兩種錯：

\boardfile[：PC 上的 Vivado xsim]{S1_csr_irq_gate.log}{2,4-7,9-15}

另外，整個 head 路徑的整合 gate（\code{tb_hd_int}）用改過的 \code{twn_csr.v} 重跑，4 個 operator、434,176 bytes 逐位元相同：

\boardfile{S2_hd_int_gate.log}{626,631,634}

\subsection{從 IP 到 device tree}

RTL 頂層 \code{twn_fuse_layer.v} 多一個 port，名字就叫 \code{interrupt}：

\snip[Verilog]{layer_irq_port}

封裝 IP 時把它宣告成 interrupt 介面。沒有這一步，Vivado 只把它當一條普通的線，PetaLinux 也不會在 device tree 寫出中斷：

\snip[tcl]{pkg_irq}

block design 把它接到 PS：

\snip[tcl]{bd_irq}

PetaLinux 從 XSA 產生的節點（開機組 E 的 \code{system.dtb}，用 \code{dtc} 反組譯）：

\boardfile[：PC 上的檢查，節點原文]{S3_binding_check.log}{24-33}

\code{interrupts = <0x00 0x59 0x04>} 的三個數字：0 表示 SPI（Shared Peripheral Interrupt，所有 CPU 共用的外部中斷），\code{0x59} = 89 是 SPI 編號，4 是 \code{IRQ_TYPE_LEVEL_HIGH}。GIC 內部的編號是 SPI 編號加 32，所以 \code{/proc/interrupts} 會顯示 121。

\subsection{driver：ISR 與 completion}

\textbf{probe 時申請中斷。}device tree 有中斷、且 \code{use_irq} 沒被關掉，就用 \code{devm_request_irq()} 註冊 ISR，再寫 IRQ 暫存器打開兩個致能；沒有中斷就留著輪詢：

\snip{drv_probe_irq}

\code{platform_get_irq_optional()} 和 \code{platform_get_irq()} 的差別是：找不到中斷時不在 kernel log 印錯誤，因為對舊 bitstream 來說這是正常的。

\textbf{ISR。}看哪一個 done 設起來，就清掉它、叫醒等它的人：

\snip{drv_isr}

\textbf{等待。}\code{wait_done()} 有中斷時睡在 completion 上，沒有時照舊輪詢：

\snip{drv_wait_done}

\textbf{啟動一幀。}和 v1.0 相比只多一行 \code{reinit_completion()}，而且位置有講究，見「原理」：

\snip{drv_bb_run}

head 的兩個函式也一樣：\code{twn_hd_start()} 在 START 之前 \code{reinit_completion(&td->hd_done)}，\code{twn_hd_wait()} 呼叫 \code{wait_done()}。

\textbf{卸載。}先關掉中斷來源，再讓 \code{devm} 釋放中斷：

\snip{drv_remove}

\textbf{板上：開機組 E。}driver 開機自動載入，說它用中斷；\code{/proc/interrupts} 有一條 GIC 121、Level：

\boardfile{E1_boot_state.log}{11-14}

\irqBoardEvidence

\subsection{原理}

\textbf{為什麼用 level，ISR 為什麼一定要清 done。}level 中斷的意思是「線拉高的期間，中斷一直成立」。GIC 在 ISR 結束（EOI）後會重新看這條線；如果 done 還沒清，線還是高，ISR 會立刻再被呼叫，形成中斷風暴（interrupt storm）。所以 ISR 一定要把中斷來源（done）清掉。level 的好處是不會漏：就算兩個 done 幾乎同時設起來，線會一直高到兩個都被處理。

\textbf{為什麼清完還要讀一次。}\code{writel()} 是 posted write：CPU 送出就繼續往下走，不等寫入真的到達 PL。ISR 若馬上返回，GIC 可能還看到舊的高電位，多進一次 ISR（這次兩個 done 都已清掉，會回傳 \code{IRQ_NONE}）。讀同一個裝置的暫存器時，Device 記憶體保證讀在寫之後到達，所以讀回來時，清除的寫入已經到了 PL。

\textbf{\code{reinit_completion()} 為什麼放在 START 之前。}completion 是一個計數加一個等待佇列：\code{complete()} 把計數加一並叫醒等待者，\code{wait_for_completion_timeout()} 在計數大於 0 時把它減一並返回。若 \code{reinit_completion()}（計數歸零）放在 START 之後，PL 若很快做完，ISR 的 \code{complete()} 可能先發生，接著被歸零，等待的人就睡到 3 秒逾時。放在 START 之前，PL 根本還沒開始，不可能有中斷搶先。

\textbf{為什麼 ISR 可以直接在硬體中斷裡做完。}它只做幾次暫存器讀寫和 \code{complete()}，都不會睡眠，可以在中斷的上下文（interrupt context）執行，不需要 threaded IRQ。

\textbf{為什麼致能要寫兩次（probe 打開、remove 關掉）。}中斷線在 FPGA 裡一直接著，driver 沒載入時致能是 reset 值 0，線不會動。probe 在 \code{devm_request_irq()} 之後才打開致能，remove 先關掉致能，再由 \code{devm} 釋放中斷；順序反過來，就可能有一個沒有人處理的中斷。

\subsection{面試會問}

\begin{qa}{你怎麼把輪詢改成中斷？從硬體講到軟體。}
硬體上，我在控制暫存器模組加了一個暫存器輸出的 level-high 中斷，等於「致能 AND done」；致能放在獨立的暫存器，因為 CTRL 每次寫入都會覆蓋原本的 bit 2，driver 一啟動就會把中斷關掉，這個錯誤我在 testbench 裡故意做出來、確認抓得到。IP 封裝時把 port 宣告成 interrupt 介面，block design 接到 PS 的 \code{pl_ps_irq0}，PetaLinux 就在 device tree 寫出 SPI 89、level high。driver 在 probe 用 \code{devm_request_irq()} 註冊 ISR；ISR 清掉 done、讀一次暫存器確保清除已經到達，再 \code{complete()}；等待的那一邊用 \code{wait_for_completion_timeout()}，\code{reinit_completion()} 放在 START 之前避免漏掉中斷。結果每幀 CPU 時間從 \pollCpuUs{} µs 降到 \irqCpuUs{} µs，fps 不變；device tree 沒有中斷時自動退回輪詢。
\end{qa}

\subsection{關鍵函式速查}

\begin{center}
\small
\begin{tabularx}{\textwidth}{l l X}
\toprule
名稱 & 位置 & 做什麼 \\
\midrule
\code{irq} 暫存器 & \code{rtl/fuse/twn_csr.v}:\lineof{csr_irq_reg} & 致能 AND done，level high \\
\code{A_IRQ} 寫入 & \code{rtl/fuse/twn_csr.v}:\lineof{csr_irq_wr} & 兩個致能 \\
interrupt 介面 & \code{vivado/package_fuse_ip.tcl}:\lineof{pkg_irq} & \code{SENSITIVITY = LEVEL_HIGH} \\
接到 PS & \code{vivado/build_bd_fuse.tcl}:\lineof{bd_irq} & \code{pl_ps_irq0} \\
\code{use_irq} & \code{twn_fuse.c}:\lineof{drv_param_irq} & 可以關掉中斷，退回輪詢 \\
\code{wait_done()} & \code{twn_fuse.c}:\lineof{drv_wait_done} & completion 或輪詢 \\
\code{twn_isr()} & \code{twn_fuse.c}:\lineof{drv_isr} & 清 done、讀回、\code{complete()} \\
probe 的中斷段 & \code{twn_fuse.c}:\lineof{drv_probe_irq} & \code{devm_request_irq()}、打開致能 \\
\bottomrule
\end{tabularx}
\end{center}
